% ==============================================================================
%  Chapter 4 -- Geometric deep learning and gauge-equivariant networks
%  Purpose: the ML half of the background, ~12-18 pages. Takes a reader from
%  "what is a neural network" to "what L-CNN does and what it cannot do",
%  ending exactly where Ch. 5 (GELT) picks up.
%  Sources: notes/papers_review.md (L-CNN sec. 1, Nagai-Tomiya, CASK).
% ==============================================================================

\section{Symmetry as an inductive bias}
\label{sec:gdl-blueprint}

\TODO{The geometric deep learning blueprint \cite{bronstein2021gdl}:
symmetry group $\to$ equivariant layers $\to$ invariant readout.
Definitions of invariance and equivariance
($f(g \cdot x) = g \cdot f(x)$); CNNs as the translation-equivariance
special case; why building the symmetry in beats data augmentation for
EXACT local symmetries (the gauge group volume grows with the lattice
volume -- augmentation cannot cover it).}

\section{Attention and transformers}
\label{sec:gdl-attention}

\TODO{Scaled dot-product attention and multi-head self-attention
\cite{vaswani2017attention}, kept minimal: queries/keys/values, softmax
scores, permutation equivariance, positional encodings (absolute bias
vs.\ rotary/RoPE \cite{su2021rope} -- both reappear as GELT variants).
The property the thesis exploits: attention weights are
CONFIGURATION-DEPENDENT, unlike convolution kernels, and they are readable
after the fact -- the hook for \cref{ch:interpretability}.}

\section{Gauge-equivariant architectures for the lattice}
\label{sec:gdl-lcnn}

\subsection{L-CNN}
\TODO{Favoni et al.\ \cite{favoni2022lcnn} in detail (this is GELT's
scaffold AND its matched-parameter baseline): parallel-transported
convolutions (L-Conv), bilinear layers (L-Bilin), trace-gated activations,
the loop-doubling universality argument (stacked L-Conv+L-Bilin generates
arbitrarily large Wilson loops). The repo implements it faithfully in
\texttt{gelt/lcnn.py} with axis-aligned transports.}

\subsection{Gauge-covariant ResNets and CASK}
\TODO{Nagai--Tomiya \cite{nagai2021covariant} (covariant smearing-like
layers, 4D non-abelian) and CASK \cite{cask2025} (gauge-covariant
transformer) -- one paragraph each; how GELT differs (invariant SCORES +
matrix-bilinear VALUES; all-shortest-path $L^1$-ball transport rather than
axis-aligned or staple-based).}

\section{What convolution cannot give you}
\label{sec:gdl-gap}

\TODO{The chapter's closing argument, setting up Ch.~5: static kernels
cannot adapt the receptive field to the configuration; physically
interesting structure (topological lumps, critical correlations) is
inhomogeneous and $\beta$-dependent. Attention is the minimal mechanism
that fixes this while keeping exact equivariance -- IF scores are gauge
invariant and values transform covariantly. That "if" is Ch.~5.}
